% ===========================================================================
%  附录 B：LQR
% ===========================================================================

\section{局部线性二次型调节器}
\label{sec:lqr}

线性二次型调节器（LQR）可在固定工作点对姿态误差、角速度和执行器代价进行系统化权衡。
令小角度状态为
$\bm x=[\bm e_R^T,\Delta\bm\omega^T]^T$，则局部模型可写为
\begin{equation}
\dot{\bm x}=\bm A\bm x+\bm B\Delta\bm u,
\qquad
\bm B=
\begin{bmatrix}\bm0\\\bm I^{-1}\bm B_M\end{bmatrix}.
\label{eq:lqr_model}
\end{equation}
小误差下 $\dot{\bm e}_R\approx\Delta\bm\omega$；非零参考角速度时还会出现坐标运输项，
不能直接沿用定点模型。

对代价函数
\begin{equation}
J=\int_0^\infty
(\bm x^T\bm Q\bm x+\Delta\bm u^T\bm R\Delta\bm u)\,\mathrm dt,
\label{eq:lqr_cost}
\end{equation}
连续代数 Riccati 方程为
\begin{equation}
\bm A^T\bm P+\bm P\bm A-
\bm P\bm B\bm R^{-1}\bm B^T\bm P+\bm Q=\bm0,
\qquad
\bm K=\bm R^{-1}\bm B^T\bm P.
\label{eq:care}
\end{equation}
控制增量为 $\Delta\bm u=-\bm K\bm x$。矩阵 $\bm Q$、$\bm R$ 的选择具有物理含义，
必须结合姿态单位、角速度单位和执行器量纲归一化；仅给出“最优”而不报告权重和工作点
不具可复现性。

冻结 $\bm B_M$ 的 LQR 只在设计点附近有效。摆角、转速、气动导数和电池状态变化会改变
实际闭环矩阵；增益调度可以在多个工作点插值，但仍需验证插值区间和共同 Lyapunov 性质。
此外，标准 LQR 不直接处理 $\pm15^\circ$ 摆角和电机转速约束，实际实现需要抗饱和、
约束分配或改用 MPC。

当前 LQR 脚本的一处分支遗漏惯量因子，使其与其他控制器的量纲不一致；
现有数值因此只用于暴露脚本问题，不进入主结论。修复后应按\cref{sec:simulation}的统一验证协议
重新生成数据。
